% =====================================================================
%  باب البسملة — الأبيات ١٠٠–١٠٧
% =====================================================================
\chapter{باب البسملة}

\textbf{البسملة}: قول «بسم الله الرحمن الرحيم»، مصدر «بَسْمَلَ» مثل: هَلَّلَ (قال لا إله إلا الله)، وحَمْدَلَ (قال الحمد لله)، وحَوْقَلَ (قال لا حول ولا قوة إلا بالله).

يتكلم الباب عن ثلاثة مواضع:
\begin{enumerate}
\item \textbf{بين السورتين} (آخر سورة مع أول التي بعدها) — وفيه خلاف القراء. \hfill (الأبيات ١٠٠–١٠٤)
\item \textbf{أول سورة براءة}. \hfill (البيت ١٠٥)
\item \textbf{الابتداء بأول السورة أو بوسطها}، وحكم الوقف على البسملة. \hfill (البيتان ١٠٦–١٠٧)
\end{enumerate}

\section{البسملة بين السورتين}

\bayt{100}{وَبَسْمَلَ بَيْنَ السُّورَتَيْنِ \rz{Nafi}{بِسُنَّةٍ}}{\rz{Kisai}{رِجَالٌ} \rz{Asim}{نَمَوْهَا} \rz{IbnKathir}{دِرْيَةً} وَتَحَمُّلَا}
\begin{mufradat}
\kalima{بسنّة}{بطريقة النبي ﷺ. والباء رمز \Qalun.}
\kalima{رجال}{الراء رمز \Kisai.}
\kalima{نَمَوها}{رفعوها ونقلوها (إلى النبي ﷺ). والنون رمز \Asim.}
\kalima{دِرية}{علمًا ومعرفة. والدال رمز \IbnKathir.}
\kalima{تحمّلا}{نقلًا عن الشيوخ (رواية).}
\end{mufradat}
\maana قرأ \Qalun{} و\Kisai{} و\Asim{} و\IbnKathir{} بإثبات البسملة بين كل سورتين، متمسكين في ذلك بسنة نقلوها عن النبي ﷺ، جامعين بين الدراية والرواية.
ومن هذه السنة: أن النبي ﷺ كان لا يعرف انقضاء السورة حتى ينزل عليه «بسم الله الرحمن الرحيم»، وكتابة الصحابة لها في المصاحف العثمانية.

\bayt{101}{وَوَصْلُكَ بَيْنَ السُّورَتَيْنِ \rz{Hamza}{فَصَاحَةٌ}}{وَصِلْ وَاسْكُتَنْ \rz{IbnAmir}{كُلٌّ} \rz{Nafi}{جَلَايَاهُ} \rz{AbuAmr}{حَصَّلَا}}
\begin{mufradat}
\kalima{فصاحة}{بيان. والفاء رمز \Hamza.}
\kalima{كلٌّ}{الكاف رمز \IbnAmir.}
\kalima{جلاياه}{الأمور الواضحة (جمع جَلِيّة). والجيم رمز \Warsh.}
\kalima{حصّلا}{الحاء رمز \AbuAmr.}
\end{mufradat}
\maana
\textbf{أولًا}: \Hamza{} يصل آخر السورة بأول التي بعدها \textbf{بدون بسملة}. وفي قوله «فصاحة» إشارة إلى حكمة الوصل: أنه يُظهر إعراب آخر السورة، كما في آخر التوبة مع أول يونس \ayah{وَهُوَ رَبُّ ٱلۡعَرۡشِ ٱلۡعَظِيمِ ۝ الٓرۚ}، ويُظهر حال همزة الوصل كآخر العاديات مع أول القارعة.

\textbf{ثانيًا}: \IbnAmir{} و\Warsh{} و\AbuAmr{} لكل منهم وجهان بين السورتين (من غير بسملة): \textbf{الوصل} و\textbf{السكت}.

\begin{qaida}[ما هو السكت بين السورتين؟]
\textbf{السكت}: قطع الصوت على آخر السورة زمنًا يسيرًا \textbf{دون تنفس}، بنية مواصلة القراءة، ثم تبدأ السورة التالية بلا بسملة. (وهو أقصر من الوقف الذي يكون معه تنفس.)
\end{qaida}

\begin{tanbih}[وصل حمزة مع سكت خلف]
\Khalaf{} عن \Hamza{} يسكت على الساكن قبل الهمز (كما سيأتي في بابه)، فإذا وصل آخر الضحى \ayah{فَحَدِّثۡ} بأول الشرح \ayah{أَلَمۡ نَشۡرَحۡ} سكت على الثاء. وهذا لا يُخرجه عن كونه \textbf{وصلًا}؛ لأن سكت خلف هذا لا يكون إلا في حال الوصل.
\end{tanbih}

\bayt{102}{وَلَا نَصَّ كَلَّا حُبَّ وَجْهٌ ذَكَرْتُهُ}{وَفِيهَا خِلَافٌ جِيدُهُ وَاضِحُ الطُّلَى}
\begin{mufradat}
\kalima{كلّا}{حرف ردع وزجر.}
\kalima{حُبّ وجهٌ}{ما أحبّ هذا الوجه! (وجهٌ: يعني الوصل والسكت).}
\kalima{جيده}{عنقه.}
\kalima{الطُّلى}{صفحات العنق، و«واضح الطلى» كناية عن الشهرة والوضوح.}
\end{mufradat}
\maana لم يرد نصٌّ عن \IbnAmir{} و\Warsh{} و\AbuAmr{} بالوصل ولا بالسكت، وإنما التخيير بين الوجهين اختيار من أهل الأداء واستحباب من شيوخ الإقراء. وكأن الناظم يزجر («كلّا») من يظن أن هناك نصًّا عنهم، ثم يمدح هذا الوجه («حُبّ وجهٌ»).
ثم قال: \textbf{وفيها} (أي في البسملة) \textbf{خلافٌ} مشهور واضح عن هؤلاء الثلاثة: فيُروى لهم أيضًا البسملة بين السورتين.

\begin{tanbih}[ليس في البيت رمز]
الكاف في «كلّا» والحاء في «حبّ» والجيم في «جيده» \textbf{ليست رموزًا}؛ لأن البيت كله في حكم الثلاثة المذكورين قبله. وهذا ما عليه المحققون.
\end{tanbih}

\begin{khulasa}[حاصل الأبيات ١٠٠–١٠٢: الأوجه بين السورتين]
\centering
\renewcommand{\arraystretch}{1.5}
\begin{tabular}{l l}
\toprule
\textbf{القراء} & \textbf{الأوجه بين السورتين}\\
\midrule
\Qalun، \IbnKathir، \Asim، \Kisai & البسملة فقط\\
\Warsh، \AbuAmr، \IbnAmir & البسملة، والوصل، والسكت\\
\Hamza & الوصل فقط (بلا بسملة)\\
\bottomrule
\end{tabular}
\end{khulasa}

\begin{tahrir}
\ithaf{9}{وَفِيهَا خِلَافٌ جِيدُهُ وَاضِحُ الطُّلَا}{وَذُو الْخُلْفِ لِلْبَصْرِيِّ وَشَامٍ تَنَقَّلَا}
\medskip
قيّد صاحب الإتحاف قول الشاطبي «وفيها خلاف» بأن وجه البسملة إنما نُقل لـ\AbiAmr{} البصري و\IbnAmir{} الشامي. ومقتضى ذلك أن \Warsh{} من طريق الشاطبية له بين السورتين الوصل والسكت فقط.
وهذه من المسائل التي يختلف فيها أهل التحرير، فخذ فيها بما يقرئك به شيخك.
\end{tahrir}

\section{البسملة بين السور الأربع الزُّهر}

\bayt{103}{وَسَكْتُهُمُ الْمُخْتَارُ دُونَ تَنَفُّسٍ}{وَبَعْضُهُمُ فِي الْأَرْبَعِ الزُّهْرِ بَسْمَلَا}
\bayt{104}{لَهُمْ دُونَ نَصٍّ وَهْوَ فِيهِنَّ سَاكِتٌ}{\qname{Hamza}{لِحَمْزَةَ} فَافْهَمْهُ وَلَيْسَ مُخَذَّلَا}
\begin{mufradat}
\kalima{سكتهم}{أي سكت الثلاثة: \Warsh{} و\AbuAmr{} و\IbnAmir.}
\kalima{المختار}{المقدَّم على الوصل.}
\kalima{الأربع الزُّهر}{أربع سور مشهورة: القيامة، والمطففين، والبلد، والهُمَزة. (الزُّهر: المشرقة، كناية عن شهرتها.)}
\kalima{مخذَّلًا}{متروك النصرة، أي ضعيفًا.}
\end{mufradat}
\maana[(البيتان ١٠٣–١٠٤)]
\textbf{أولًا}: السكت للثلاثة (\Warsh{} و\AbuAmr{} و\IbnAmir) هو الوجه المختار المقدَّم على الوصل، لأن فيه تنبيهًا على انتهاء السورة. ويكون بلا تنفس.

\textbf{ثانيًا}: بعض أهل الأداء استحب للقراء مذهبًا خاصًا في أربع سور هي: \textbf{القيامة، المطففين، البلد، الهمزة}، وذلك لبشاعة اللفظ عند وصلها بما قبلها (مثل وصل \ayah{وَأَهۡلُ ٱلۡمَغۡفِرَةِ} بـ\ayah{لَآ أُقۡسِمُ}، أو \ayah{وَٱلۡأَمۡرُ يَوۡمَئِذࣲ لِّلَّهِ} بـ\ayah{وَيۡلࣱ لِّلۡمُطَفِّفِينَ}):
\begin{itemize}
\item من كان يسكت في غيرها (\Warsh{} و\AbuAmr{} و\IbnAmir) \ila{} \textbf{يبسمل} بين هذه الأربع وما قبلها.
\item من كان يصل في غيرها (هؤلاء الثلاثة في وجه الوصل، و\Hamza) \ila{} \textbf{يسكت} بين هذه الأربع وما قبلها.
\end{itemize}
وهذا كله بلا نص عن القراء، بل هو اختيار. والناظم يقول: «فافهمه وليس مخذَّلا» أي هذا المذهب ليس ضعيفًا. لكن \textbf{المحققون على عدم التفريق} بين هذه السور وغيرها، وهو المذهب المختار الذي عليه العمل.

\begin{amthila}[المواضع الأربعة]
\begin{itemize}
\item بين المدثر و\textbf{القيامة}: \ayah{هُوَ أَهۡلُ ٱلتَّقۡوَىٰ وَأَهۡلُ ٱلۡمَغۡفِرَةِ ۝ لَآ أُقۡسِمُ بِيَوۡمِ ٱلۡقِيَٰمَةِ}.
\item بين الانفطار و\textbf{المطففين}: \ayah{وَٱلۡأَمۡرُ يَوۡمَئِذࣲ لِّلَّهِ ۝ وَيۡلࣱ لِّلۡمُطَفِّفِينَ}.
\item بين الفجر و\textbf{البلد}: \ayah{وَٱدۡخُلِي جَنَّتِي ۝ لَآ أُقۡسِمُ بِهَٰذَا ٱلۡبَلَدِ}.
\item بين العصر و\textbf{الهمزة}: \ayah{وَتَوَاصَوۡاْ بِٱلصَّبۡرِ ۝ وَيۡلࣱ لِّكُلِّ هُمَزَةࣲ}.
\end{itemize}
\end{amthila}

\begin{tahrir}
\ithaf{10}{وَبَسْمِلْ بِزُهْرٍ إِنْ تُبَسْمِلْ بِغَيْرِهَا}{وَإِنْ تَسْكُتُ اسْكُتْ بَعْدَ مَا أَنْ تُبَسْمِلَا}
\ithaf{11}{وَإِنْ تَصِلَنْ فَاسْكُتْ بِهَا ثُمَّ صِلْ وَإِنْ}{بَدَأْتَ بِهَا بَسْمِلْ بِهَا وَبِمَا تَلَا}
\ithaf{12}{فَبَسْمِلْ كَذَا اسْكُتْ ثُمَّ إِنْ تَسْكُتَا بِهَا}{فَفِي غَيْرِهَا اسْكُتْ صِلْ وَإِنْ تَصِلًا صِلَا}
\medskip
\textbf{خلاصة ما يتضح من هذه الأبيات} (لمن أخذ بمذهب التفريق في الزُّهر):
\begin{itemize}
\item إن كنت تبسمل في غير الزهر \ila{} فبسمل في الزهر أيضًا.
\item إن كنت تسكت في غيرها \ila{} ففي الزهر البسملة، ويجوز السكت أيضًا.
\item إن كنت تصل في غيرها \ila{} ففي الزهر السكت، ويجوز الوصل أيضًا.
\end{itemize}
وفي بقية البيتين ١١–١٢ تفصيلات دقيقة في ترتيب الأوجه عند الجمع، تُتلقّى من الشيخ عند الأداء.
\end{tahrir}

\section{أول سورة براءة}

\bayt{105}{وَمَهْمَا تَصِلْهَا أَوْ بَدَأْتَ بَرَاءَةً}{لِتَنْزِيلِهَا بِالسَّيْفِ لَسْتَ مُبَسْمِلَا}
\begin{mufradat}
\kalima{تصلها}{تصل سورة براءة (التوبة) بالأنفال قبلها.}
\kalima{بالسيف}{بالقتال والتهديد ونبذ العهود.}
\end{mufradat}
\maana سورة براءة (التوبة) \textbf{لا بسملة في أولها لأحد من القراء}، سواء وصلتها بالأنفال قبلها أو ابتدأت القراءة بها. والعلة: أنها نزلت بالسيف، أي بالأمر بالقتال ونبذ العهود والوعيد. وقد نُقل عن علي رضي الله عنه: «بسم الله أمان، وبراءة ليس فيها أمان؛ نزلت بالسيف».

\begin{tahrir}
\ithaf{13}{وَلِلْكُلِّ قِفْ صِلْ فِي عَلِيمٍ بَرَاءَةٌ}{أَوِ اسْكُتْ وَبَيْنَ النَّاسِ وَالْحَمْدِ بَسْمِلَا}
\medskip
\begin{itemize}
\item بين آخر الأنفال \ayah{إِنَّ ٱللَّهَ بِكُلِّ شَيۡءٍ عَلِيمُۢ} وأول براءة: لكل القراء ثلاثة أوجه: \textbf{الوقف، والسكت، والوصل} (كلها بلا بسملة).
\item بين آخر سورة الناس وأول الفاتحة: \textbf{البسملة لكل القراء} (حتى من مذهبه الوصل أو السكت)؛ لأن الفاتحة مبتدأ بها حكمًا، إذ ليس قبلها شيء في الحقيقة.
\end{itemize}
\end{tahrir}

\section{الابتداء بأول السورة وبوسطها}

\bayt{106}{وَلَا بُدَّ مِنْهَا فِي ابْتِدَائِكَ سُورَةً}{سِوَاهَا وَفِي الْأَجْزَاءِ خُيِّرَ مَنْ تَلَا}
\begin{mufradat}
\kalima{منها}{من البسملة.}
\kalima{سواها}{غير براءة.}
\kalima{الأجزاء}{ما بعد أول السورة ولو بآية أو كلمة (أواسط السور).}
\kalima{خُيِّر من تلا}{خُيِّر القارئ.}
\end{mufradat}
\maana
\begin{itemize}
\item إذا \textbf{ابتدأت} القراءة بأول أي سورة غير براءة \ila{} \textbf{لا بد} من البسملة لجميع القراء، حتى \Hamza{} ومن له السكت أو الوصل؛ لأن مذاهبهم تلك إنما هي \textbf{بين} السورتين.
\item إذا ابتدأت من \textbf{وسط السورة} (ولو من الآية الثانية) \ila{} القارئ مخيّر بين البسملة وتركها، لكل القراء.
\end{itemize}

\begin{tanbih}[الابتداء بوسط براءة]
لا فرق في هذا التخيير بين أجزاء براءة وأجزاء غيرها على الصحيح، وإن منع بعضهم البسملة في أجزاء براءة إلحاقًا لها بأولها.
\end{tanbih}

\bayt{107}{وَمَهْمَا تَصِلْهَا مَعْ أَوَاخِرِ سُورَةٍ}{فَلَا تَقِفَنَّ الدَّهْرَ فِيهَا فَتَثْقُلَا}
\begin{mufradat}
\kalima{تصلها}{تصل البسملة.}
\kalima{فيها}{عليها.}
\kalima{تثقلا}{تصير مستثقَلًا عند أئمة القراءة.}
\end{mufradat}
\maana إذا وصلت البسملة بآخر السورة فلا تقف عليها أبدًا، بل يجب وصلها بأول السورة التالية؛ لأن البسملة شُرعت لأول السورة لا لآخرها، فإذا وقفت عليها بعد وصلها بآخر السورة أوهمت أنها من آخر السورة السابقة.

\begin{khulasa}[الأوجه بين السورتين لمن يبسمل]
\centering
\renewcommand{\arraystretch}{1.5}
\begin{tabular}{c l c}
\toprule
\textbf{الوجه} & \textbf{الصورة} & \textbf{الحكم}\\
\midrule
١ & قطع الجميع: آخر السورة \textbf{(وقف)} البسملة \textbf{(وقف)} أول السورة & جائز\\
٢ & قطع الأول ووصل الثاني بالثالث & جائز\\
٣ & وصل الجميع & جائز\\
٤ & وصل آخر السورة بالبسملة ثم الوقف عليها & \textbf{ممنوع}\\
\bottomrule
\end{tabular}
\end{khulasa}

\section{خلاصة الباب}
\begin{khulasa}[ما لكل قارئ بين كل سورتين (ما عدا الأنفال وبراءة)]
\centering
\renewcommand{\arraystretch}{1.5}
\begin{tabular}{l c l}
\toprule
\textbf{القراء} & \textbf{عدد الأوجه} & \textbf{تفصيلها}\\
\midrule
\Qalun، \IbnKathir، \Asim، \Kisai & ٣ & أوجه البسملة الثلاثة الجائزة\\
\Warsh، \AbuAmr، \IbnAmir & ٥ & أوجه البسملة الثلاثة + الوصل + السكت\\
\Hamza & ١ & الوصل بلا بسملة\\
\midrule
الجميع: بين الأنفال وبراءة & ٣ & الوقف، والسكت، والوصل (بلا بسملة)\\
الجميع: بين الناس والفاتحة & — & البسملة\\
الجميع: الابتداء بأول سورة & — & البسملة (إلا براءة)\\
الجميع: الابتداء بوسط سورة & — & التخيير\\
\bottomrule
\end{tabular}
\end{khulasa}
